%!TeX program = xelatex
\documentclass{SYSUReport}
\usepackage{xeCJK}
\usepackage{tikz}
\usepackage{graphicx}
\usepackage{makecell}
\usepackage{float}
\usepackage{listings}
\usepackage{xcolor}
\usepackage{multirow}
\usepackage{booktabs}
\usepackage{amsmath}
\usepackage{amssymb}

\lstset{
    basicstyle=\ttfamily\small,
    breaklines=true,
    columns=flexible,
    frame=single,
    numbers=left,
    numberstyle=\tiny\color{gray},
    keywordstyle=\color{blue},
    commentstyle=\color{green!50!black},
    stringstyle=\color{red},
    showstringspaces=false
}

\setCJKmainfont{simsun.ttc}

\headl{中山大学}
\headc{VLSI 划分实验}
\headr{实验1}
\lessonTitle{VLSI 图划分实验}
\reportTitle{超图二划分算法：FM算法与模拟退火算法}
\stuname{钟旺烜、赵施琦、郑晋凯}
\stuid{23336342、23336324、23336329}
\inst{计算机学院}
\major{计算机科学与技术}
\teacher{沈明华}
\date{2026年5月17日}

\begin{document}
\cover
\thispagestyle{empty}
\clearpage

\pagenumbering{Roman}
\setcounter{page}{1}
\tableofcontents
\clearpage

\pagenumbering{arabic}
\setcounter{page}{1}

\section{实验目的}
\subsection{题目要求}
本实验要求实现超图二划分算法，将给定的超图划分为两个大小均衡的子集，使得跨越两个子集的超边数（cut size）最小化。

\textbf{具体要求}：
\begin{itemize}
\item 实现FM（Fiduccia-Mattheyses）算法进行超图二划分
\item 实现模拟退火算法进行超图二划分
\item 在ISPD98基准测试集（ibm01--ibm18）上验证算法正确性与性能
\item 划分平衡约束：两个子集的大小比例在$50\% \pm \epsilon$范围内，其中$\epsilon = 0.02$
\item 评估指标为cut size，即跨越两个划分的超边数量
\end{itemize}

\subsection{实验目的}
通过实现超图二划分算法，深入理解VLSI物理设计中划分问题的基本原理和求解方法。掌握FM算法的增益计算、增量更新与多起点策略。理解模拟退火算法在组合优化问题中的应用，包括邻域结构设计、温度调度与接受准则。通过在ISPD98基准测试集上的实验，分析两种算法的划分质量与性能特征。

\section{实验过程}

\subsection{超图划分问题定义}
超图$H = (V, E)$是图的推广，其中$V$为顶点集，$E$为超边集。每条超边$e \in E$可以连接两个或更多顶点。超图二划分问题将顶点集$V$划分为两个不相交的子集$X$和$Y$（$X \cup Y = V$，$X \cap Y = \emptyset$），使得跨越两个子集的超边数最小：

\begin{equation}
\text{cut}(X, Y) = |\{e \in E \mid e \cap X \neq \emptyset \text{ 且 } e \cap Y \neq \emptyset\}|
\end{equation}

同时满足平衡约束：
\begin{equation}
\frac{|X|}{|V|} \in [0.5 - \epsilon, 0.5 + \epsilon]
\end{equation}

其中$\epsilon = 0.02$，即每个子集的大小占总顶点数的48\%--52\%。

在VLSI设计中，超图的顶点对应电路模块，超边对应信号线网（net）。划分问题旨在将电路分为两部分，使得两部分之间的连线数最少，从而减少芯片间的通信开销。

\subsection{数据结构设计}
\subsubsection{超图表示}
采用面向对象的方式表示超图，核心数据结构包括：

\begin{lstlisting}[language=C++, caption=核心数据结构定义]
class Node {
    int index;
    vector<Net*> nets;
public:
    void add_net(Net *net);
    int get_index() { return this->index; }
    vector<Net*> &get_nets() { return this->nets; }
};

class Net {
    int index;
    vector<Node*> nodes;
public:
    void add_node(Node *node);
    int get_index() { return this->index; }
    vector<Node*> &get_nodes() { return this->nodes; }
};

class Graph {
    vector<Node*> nodes;
    vector<Net*> nets;
    map<int, Node*> node_map;
    map<int, Net*> net_map;
public:
    Node *get_or_create_node(int index);
    Net *add_net(int index);
    vector<Node*> &get_nodes() { return this->nodes; }
    vector<Net*> &get_nets() { return this->nets; }
};
\end{lstlisting}

\textbf{设计要点}：
\begin{itemize}
\item Node与Net之间通过指针双向关联，支持从顶点快速访问其关联的超边，也支持从超边快速访问其关联的顶点
\item Graph使用\texttt{node\_map}和\texttt{net\_map}实现按索引快速查找，避免线性搜索
\item \texttt{get\_or\_create\_node}方法在节点不存在时自动创建，简化了基准文件的读取逻辑
\end{itemize}

\subsubsection{基准文件格式}
ISPD98基准文件采用hMETIS格式：

\begin{lstlisting}[caption=hMETIS超图文件格式]
<超边数> <顶点数>
<超边1连接的顶点列表>
<超边2连接的顶点列表>
...
\end{lstlisting}

读取时，每行对应一条超边，行中的整数表示该超边连接的顶点编号。

\subsection{FM算法实现}
\subsubsection{FM算法原理}
Fiduccia-Mattheyses算法是一种经典的超图划分启发式算法，基于Kernighan-Lin算法改进而来。其核心思想是通过迭代地移动顶点来逐步改善划分质量。与KL算法不同，FM算法每次只移动一个顶点，且使用增益（gain）来指导顶点的选择。

\textbf{增益定义}：顶点$v$的增益$g(v)$定义为将$v$从当前子集移动到另一个子集时，cut size的减少量：

\begin{equation}
g(v) = |\{e \in \text{nets}(v) \mid e \text{ 将从cut中移除}\}| - |\{e \in \text{nets}(v) \mid e \text{ 将被加入cut}\}|
\end{equation}

具体而言，对于顶点$v$所在的每条超边$e$：
\begin{itemize}
\item 若$e$在$v$的当前子集中仅有$v$一个顶点（critical net），移动$v$将使$e$从cut中移除，增益$+1$
\item 若$e$在$v$的目标子集中没有顶点，移动$v$将使$e$被加入cut，增益$-1$
\end{itemize}

\subsubsection{增益计算}
增益计算函数的实现如下：

\begin{lstlisting}[language=C++, caption=增益计算函数]
int Solution::calculate_gain(Node *node, map<int, int> &partition_map) {
    int gain = 0;
    int node_part = partition_map[node->get_index()];

    for(Net *net : node->get_nets()) {
        int count_in_X = 0, count_in_Y = 0;
        for(Node *n : net->get_nodes()) {
            if(partition_map[n->get_index()] == 0) count_in_X++;
            else count_in_Y++;
        }

        if(node_part == 0) {
            if(count_in_X == 1 && count_in_Y > 0) gain++;
            if(count_in_Y == 0) gain--;
        } else {
            if(count_in_Y == 1 && count_in_X > 0) gain++;
            if(count_in_X == 0) gain--;
        }
    }
    return gain;
}
\end{lstlisting}

\subsubsection{FM Pass}
一次FM Pass是FM算法的核心迭代过程：

\begin{enumerate}
\item 计算所有未锁定顶点的增益
\item 选择满足平衡约束且增益最大的未锁定顶点
\item 移动该顶点到另一个子集，并锁定该顶点
\item 增量更新受影响顶点的增益
\item 重复步骤2--4直到所有顶点被锁定或无法移动
\item 回溯到累积增益最大的划分状态
\end{enumerate}

\begin{lstlisting}[language=C++, caption=FM Pass核心逻辑]
void Solution::fm_pass(Graph &graph, map<int, int> &partition_map,
                       int &size_X, int &size_Y) {
    map<int, int> gains;
    vector<bool> locked(total_size + 1, false);
    int cumulative_gain = 0, best_cumulative = 0;
    map<int, int> best_partition = partition_map;

    for(Node *node : nodes)
        gains[node->get_index()] = calculate_gain(node, partition_map);

    while(locked_count < total_size) {
        int best_gain = INT_MIN, best_node_idx = -1;
        for(Node *node : nodes) {
            int idx = node->get_index();
            if(!locked[idx] && is_balanced(...)) {
                if(gains[idx] > best_gain) {
                    best_gain = gains[idx];
                    best_node_idx = idx;
                }
            }
        }
        if(best_node_idx == -1) break;

        partition_map[best_node_idx] = 1 - old_part;
        locked[best_node_idx] = true;
        cumulative_gain += best_gain;

        if(cumulative_gain > best_cumulative) {
            best_cumulative = cumulative_gain;
            best_partition = partition_map;
        }
        update_gains_incremental(best_node, partition_map, gains, locked);
    }
    partition_map = best_partition;
}
\end{lstlisting}

\textbf{关键设计}：
\begin{itemize}
\item \textbf{锁定机制}：已移动的顶点在当前pass中不再移动，避免循环
\item \textbf{累积增益回溯}：即使中间步骤增益为负，仍可能后续获得更大改善，因此记录累积增益最大的状态
\item \textbf{平衡约束检查}：每次移动前验证移动后的划分是否满足平衡约束
\end{itemize}

\subsubsection{增量增益更新}
移动一个顶点后，只需更新与其共享超边的顶点的增益，而非重新计算所有顶点的增益。这是FM算法高效性的关键：

\begin{lstlisting}[language=C++, caption=增量增益更新]
void Solution::update_gains_incremental(Node *moved_node,
    map<int, int> &partition_map, map<int, int> &gains,
    vector<bool> &locked) {
    int moved_part = partition_map[moved_node->get_index()];
    int from_part = 1 - moved_part;

    for(Net *net : moved_node->get_nets()) {
        int count_in_from = 0, count_in_to = 0;
        Node* critical_node_from = nullptr;
        Node* critical_node_to = nullptr;

        for(Node *n : net->get_nodes()) {
            if(partition_map[n->get_index()] == from_part) {
                count_in_from++;
                if(count_in_from == 1) critical_node_from = n;
            } else {
                count_in_to++;
                if(count_in_to == 1) critical_node_to = n;
            }
        }

        if(count_in_from == 0) {
            for(Node *n : net->get_nodes())
                if(!locked[n->get_index()] && partition_map[...] == moved_part)
                    gains[n->get_index()]--;
        } else if(count_in_from == 1) {
            if(critical_node_from && !locked[critical_node_from->get_index()])
                gains[critical_node_from->get_index()]++;
        }

        if(count_in_to == 0) {
            for(Node *n : net->get_nodes())
                if(!locked[n->get_index()] && partition_map[...] == from_part)
                    gains[n->get_index()]++;
        } else if(count_in_to == 1) {
            if(critical_node_to && !locked[critical_node_to->get_index()])
                gains[critical_node_to->get_index()]--;
        }
    }
}
\end{lstlisting}

增量更新的核心逻辑基于超边的临界状态变化：
\begin{itemize}
\item \textbf{from侧顶点数为0}：移动后，该超边在moved\_part侧的所有未锁定顶点增益$-1$（因为该超边不再被cut，移动它们不会减少cut）
\item \textbf{from侧顶点数为1}：该唯一顶点成为临界顶点，增益$+1$（移动它可使超边从cut中移除）
\item \textbf{to侧顶点数为0}：from侧所有未锁定顶点增益$+1$（移动后超边将不再被cut）
\item \textbf{to侧顶点数为1}：该唯一顶点失去临界性，增益$-1$（移动它不再使超边从cut中移除）
\end{itemize}

\subsubsection{初始化策略}
FM算法对初始划分敏感，本实验实现了两种初始化策略：

\textbf{1. 随机初始化}（\texttt{initialize\_partition}）：
\begin{itemize}
\item 使用Mersenne Twister随机数生成器打乱顶点顺序
\item 将前一半顶点分配到子集$X$，后一半分配到子集$Y$
\item 不同运行使用不同随机种子
\end{itemize}

\textbf{2. BFS初始化}（\texttt{initialize\_partition\_bfs}）：
\begin{itemize}
\item 随机选择起始顶点，通过BFS遍历超图
\item 将BFS先访问到的一半顶点分配到子集$X$，其余分配到$Y$
\item BFS初始化倾向于将局部连通的顶点分配到同一子集，可能产生更好的初始划分
\end{itemize}

\subsubsection{多起点策略与OpenMP并行化}
由于FM算法是局部搜索算法，容易陷入局部最优。本实验采用多起点策略：

\begin{lstlisting}[language=C++, caption=多起点FM划分]
#define MULTI_START 20
#define MAX_PASSES 50

void Solution::fm_partition(Graph &graph, set<int> &X, set<int> &Y) {
    #pragma omp parallel
    {
        map<int, int> local_partition;
        int local_best_cut = INT_MAX;
        map<int, int> local_best_partition;

        #pragma omp for schedule(dynamic)
        for(int run = 0; run < MULTI_START; run++) {
            unsigned int seed = 42 + run * 1000;
            if(run % 2 == 0)
                initialize_partition(graph, local_partition, ...);
            else
                initialize_partition_bfs(graph, local_partition, ...);

            for(int pass = 0; pass < MAX_PASSES; pass++) {
                int prev_cut = calculate_cut_size(graph, local_partition);
                fm_pass(graph, local_partition, ...);
                int new_cut = calculate_cut_size(graph, local_partition);

                if(new_cut < local_best_cut) {
                    local_best_cut = new_cut;
                    local_best_partition = local_partition;
                }
                if(new_cut >= prev_cut && pass > 5) break;
            }
        }

        #pragma omp critical
        {
            if(local_best_cut < best_cut) {
                best_cut = local_best_cut;
                best_partition = local_best_partition;
            }
        }
    }
}
\end{lstlisting}

\textbf{关键设计}：
\begin{itemize}
\item 共执行20次独立运行（\texttt{MULTI\_START = 20}），偶数次运行使用随机初始化，奇数次使用BFS初始化
\item 每次运行最多执行50个FM Pass（\texttt{MAX\_PASSES = 50}），若连续无改善且已执行超过5个pass则提前终止
\item 使用OpenMP并行化多起点运行，\texttt{schedule(dynamic)}实现动态负载均衡
\item 每个线程维护局部最优解，最后通过\texttt{\#pragma omp critical}合并全局最优解
\end{itemize}

\subsection{模拟退火算法实现}
\subsubsection{模拟退火算法原理}
模拟退火（Simulated Annealing, SA）算法是一种基于物理退火过程的全局优化启发式算法。其核心思想是在搜索过程中以一定概率接受劣解，从而跳出局部最优。

算法的关键参数：
\begin{itemize}
\item 初始温度$T_0 = 10000$，最终温度$T_f = 0.001$
\item 冷却速率$\alpha = 0.995$，即$T_{k+1} = \alpha \cdot T_k$
\item 每个温度下的迭代次数$\text{ITER} = 100$
\item 接受准则：对于目标函数变化$\Delta$，若$\Delta < 0$则接受，否则以概率$e^{-\Delta/T}$接受
\end{itemize}

\subsubsection{邻域结构设计}
本实验设计了两种邻域操作：

\textbf{1. 单顶点移动}（概率0.5）：
\begin{itemize}
\item 随机选择一个顶点，将其从当前子集移动到另一个子集
\item 移动前检查平衡约束
\item 计算移动对cut size的增量变化$\delta$
\end{itemize}

\textbf{2. 顶点交换}（概率0.5）：
\begin{itemize}
\item 随机选择两个位于不同子集的顶点，交换它们的划分
\item 交换操作不改变两个子集的大小，天然满足平衡约束
\item 计算交换对cut size的总增量$\delta_1 + \delta_2$
\end{itemize}

\begin{lstlisting}[language=C++, caption=模拟退火核心逻辑]
void Solution::sa_partition(Graph &graph, set<int> &X, set<int> &Y) {
    double temperature = INITIAL_TEMP;

    while(temperature > FINAL_TEMP) {
        for(int iter = 0; iter < ITERATIONS_PER_TEMP; iter++) {
            if(dis(gen) < 0.5) {
                // 单顶点移动
                int node_id = node_ids[node_dist(gen)];
                if(!is_balanced(new_size_X, new_size_Y, total_size))
                    continue;
                int delta = calculate_delta(graph, partition_map, node_id);
                partition_map[node_id] = 1 - partition_map[node_id];
                if(delta < 0 || dis(gen) < exp(-delta / temperature)) {
                    current_cut += delta;
                    if(current_cut < best_cut) {
                        best_cut = current_cut;
                        best_partition = partition_map;
                    }
                } else {
                    partition_map[node_id] = 1 - partition_map[node_id];
                }
            } else {
                // 顶点交换
                int node1 = node_ids[node_dist(gen)];
                int node2 = node_ids[node_dist(gen)];
                if(partition_map[node1] == partition_map[node2]) continue;
                int delta1 = calculate_delta(..., node1);
                partition_map[node1] = 1 - partition_map[node1];
                int delta2 = calculate_delta(..., node2);
                partition_map[node2] = 1 - partition_map[node2];
                int total_delta = delta1 + delta2;
                if(total_delta < 0 || dis(gen) < exp(-total_delta / temperature)) {
                    current_cut += total_delta;
                } else {
                    partition_map[node1] = 1 - partition_map[node1];
                    partition_map[node2] = 1 - partition_map[node2];
                }
            }
        }
        temperature *= COOLING_RATE;
    }
}
\end{lstlisting}

\subsubsection{增量变化计算}
模拟退火算法中，\texttt{calculate\_delta}函数计算移动单个顶点对cut size的增量变化：

\begin{lstlisting}[language=C++, caption=增量变化计算]
int Solution::calculate_delta(Graph &graph, map<int, int> &partition_map,
                             int node_id) {
    int delta = 0;
    int part = partition_map[node_id];
    Node *node = ...;

    for(Net *net : node->get_nets()) {
        int count_X = 0, count_Y = 0;
        for(Node *n : net->get_nodes()) {
            if(partition_map[n->get_index()] == 0) count_X++;
            else count_Y++;
        }

        if(part == 0) {
            if(count_X == 1 && count_Y > 0) delta--;
            else if(count_Y == 0) delta++;
        } else {
            if(count_Y == 1 && count_X > 0) delta--;
            else if(count_X == 0) delta++;
        }
    }
    return delta;
}
\end{lstlisting}

注意：\texttt{calculate\_delta}返回的是移动顶点后cut size的\textbf{变化量}（负值表示改善），而FM算法中\texttt{calculate\_gain}返回的是\textbf{增益}（正值表示改善），两者符号相反。

\subsection{平衡约束}
两种算法均使用相同的平衡约束检查，要求移动后两个子集的大小比例在$[0.5 - \epsilon, 0.5 + \epsilon]$范围内：

\begin{lstlisting}[language=C++, caption=平衡约束检查]
#define EPSILON 0.02

bool Solution::is_balanced(int size_X, int size_Y,
                           int total_size, int moving_from) {
    int new_size_X = size_X, new_size_Y = size_Y;
    if(moving_from == 0) { new_size_X--; new_size_Y++; }
    else { new_size_X++; new_size_Y--; }
    double ratio_X = (double)new_size_X / total_size;
    return (ratio_X >= 0.5 - EPSILON && ratio_X <= 0.5 + EPSILON);
}
\end{lstlisting}

\subsection{自动化测试脚本}
编写Shell脚本对ISPD98全部18个基准测试进行批量测试：

\begin{lstlisting}[language=bash, caption=自动化测试脚本]
BENCHMARK_DIR="ISPD_benchmark"
RESULT_FILE="benchmark_results.txt"

for i in 01 02 03 04 05 06 07 08 09 10 11 12 13 14 15 16 17 18; do
    benchmark_file="${BENCHMARK_DIR}/ibm${i}.hgr"
    if [ -f "$benchmark_file" ]; then
        output=$(./main "$benchmark_file" 2>&1)
        cut_size=$(echo "$output" | grep "Cut size:" | awk '{print $3}')
        ratio_x=$(echo "$output" | grep "Ratio X:" | awk '{print $3}')
        ratio_y=$(echo "$output" | grep "Ratio Y:" | awk '{print $3}')
        printf "%-15s %-15s X: %-6s Y: %-6s\n" \
            "ibm${i}.hgr" "$cut_size" "$ratio_x" "$ratio_y"
    fi
done
\end{lstlisting}

\subsection{实验环境}
实验使用的CPU为锐龙R9 7940H，具有4个物理核心，频率为4.00GHz。编译器为g++，编译选项为\texttt{-std=c++17 -fopenmp -O2}。FM算法使用OpenMP进行多起点并行化。

\section{实验结果}

\subsection{ISPD98基准测试集信息}
ISPD98基准测试集包含18个电路，规模从约1.2万到21万顶点不等：

\begin{table}[H]
\centering
\caption{ISPD98基准测试集规模}
\begin{tabular}{|c|c|c|}
\hline
基准 & 顶点数 & 超边数 \\
\hline
ibm01 & 12752 & 14111 \\
\hline
ibm02 & 19601 & 19584 \\
\hline
ibm03 & 23136 & 27401 \\
\hline
ibm04 & 27507 & 31970 \\
\hline
ibm05 & 29347 & 28446 \\
\hline
ibm06 & 32498 & 34826 \\
\hline
ibm07 & 45926 & 48117 \\
\hline
ibm08 & 51309 & 50513 \\
\hline
ibm09 & 53395 & 60902 \\
\hline
ibm10 & 69429 & 75196 \\
\hline
ibm11 & 70558 & 81454 \\
\hline
ibm12 & 71076 & 77240 \\
\hline
ibm13 & 84199 & 99666 \\
\hline
ibm14 & 147605 & 152772 \\
\hline
ibm15 & 161570 & 186608 \\
\hline
ibm16 & 183484 & 190048 \\
\hline
ibm17 & 185495 & 189581 \\
\hline
ibm18 & 210613 & 201920 \\
\hline
\end{tabular}
\end{table}

\subsection{FM算法划分结果}
FM算法（20次多起点，每次最多50个pass）在ISPD98基准测试集上的划分结果：

\begin{table}[H]
\centering
\caption{FM算法划分结果}
\begin{tabular}{|c|c|c|c|}
\hline
基准 & Cut Size & $|X|/|V|$ & $|Y|/|V|$ \\
\hline
ibm01 & 517 & 0.5124 & 0.4876 \\
\hline
ibm02 & 1116 & 0.5199 & 0.4801 \\
\hline
ibm03 & 2104 & 0.5191 & 0.4809 \\
\hline
ibm04 & 1424 & 0.5141 & 0.4859 \\
\hline
ibm05 & 3616 & 0.5200 & 0.4800 \\
\hline
ibm06 & 2581 & 0.5200 & 0.4800 \\
\hline
ibm07 & 2235 & 0.5186 & 0.4814 \\
\hline
ibm08 & 3601 & 0.5175 & 0.4825 \\
\hline
ibm09 & 3825 & 0.5099 & 0.4901 \\
\hline
ibm10 & 4134 & 0.5069 & 0.4931 \\
\hline
ibm11 & 5051 & 0.5158 & 0.4842 \\
\hline
ibm12 & 5194 & 0.5191 & 0.4809 \\
\hline
ibm13 & 4116 & 0.5193 & 0.4807 \\
\hline
ibm14 & 8546 & 0.4903 & 0.5097 \\
\hline
ibm15 & 11405 & 0.5200 & 0.4800 \\
\hline
ibm16 & 7781 & 0.5107 & 0.4893 \\
\hline
ibm17 & 9389 & 0.5200 & 0.4800 \\
\hline
ibm18 & 14703 & 0.5200 & 0.4800 \\
\hline
\end{tabular}
\end{table}

所有划分结果均满足平衡约束（$|X|/|V| \in [0.48, 0.52]$），验证了算法的正确性。

\subsection{与已知最优解的对比}
将FM算法的结果与当前已知最优解（来自SpecPart/hMETIS/KaHyPar等先进划分器）进行对比。下表列出了不同平衡约束$\epsilon$下的已知最优cut size：

\begin{table}[H]
\centering
\caption{FM算法与已知最优解对比（$\epsilon = 2\%$）}
\begin{tabular}{|c|c|c|c|c|c|c|}
\hline
\multirow{2}{*}{基准} & \multirow{2}{*}{FM Cut Size} & \multicolumn{4}{c|}{已知最优（不同$\epsilon$）} & \multirow{2}{*}{差距($\epsilon=2\%$)} \\
\cline{3-6}
 & & $\epsilon=1\%$ & $\epsilon=2\%$ & $\epsilon=5\%$ & $\epsilon=10\%$ & \\
\hline
ibm01 & 517 & 203 & 200 & 180 & 166 & 158.5\% \\
\hline
ibm02 & 1116 & 341 & 307 & 262 & 262 & 263.5\% \\
\hline
ibm03 & 2104 & 954 & 951 & 950 & 950 & 121.3\% \\
\hline
ibm04 & 1424 & 580 & 573 & 514 & 388 & 148.5\% \\
\hline
ibm05 & 3616 & 1719 & 1706 & 1697 & 1645 & 111.8\% \\
\hline
ibm06 & 2581 & 976 & 962 & 871 & 728 & 168.3\% \\
\hline
ibm07 & 2235 & 910 & 878 & 818 & 764 & 154.6\% \\
\hline
ibm08 & 3601 & 1140 & 1140 & 1140 & 1120 & 215.9\% \\
\hline
ibm09 & 3825 & 625 & 620 & 620 & 519 & 516.9\% \\
\hline
ibm10 & 4134 & 1285 & 1253 & 1244 & 1244 & 229.9\% \\
\hline
ibm11 & 5051 & 1065 & 1051 & 951 & 763 & 380.5\% \\
\hline
ibm12 & 5194 & 1934 & 1919 & 1871 & 1841 & 170.7\% \\
\hline
ibm13 & 4116 & 837 & 831 & 831 & 655 & 395.5\% \\
\hline
ibm14 & 8546 & 1842 & 1842 & 1794 & 1509 & 363.7\% \\
\hline
ibm15 & 11405 & 2743 & 2730 & 2530 & 2135 & 317.8\% \\
\hline
ibm16 & 7781 & 1975 & 1827 & 1695 & 1619 & 326.0\% \\
\hline
ibm17 & 9389 & 2314 & 2270 & 2186 & 1989 & 313.6\% \\
\hline
ibm18 & 14703 & 1521 & 1521 & 1521 & 1520 & 866.7\% \\
\hline
\end{tabular}
\end{table}

从对比中可以看出，基本的FM算法与已知最优解之间存在较大差距。这是因为已知最优解通常由多层划分（multilevel partitioning）算法（如hMETIS、KaHyPar、SpecPart等）获得，这些算法通过粗化-初始划分-细化三个阶段，在粗化后的较小超图上搜索，再逐层细化回原始超图，能够更有效地跳出局部最优。同时，随着平衡约束$\epsilon$的放宽，已知最优解的cut size显著降低，说明更大的平衡容忍度允许算法找到更优的划分。

\begin{figure}[H]
\centering
\includegraphics[width=0.85\textwidth]{figures/Figure_1.png}
\caption{差距可视化}
\end{figure}

\begin{figure}[H]
\centering
\includegraphics[width=0.85\textwidth]{figures/Figure_2.png}
\caption{多维度对比可视化}
\end{figure}

\section{加分项}

本实验在满足基础实验要求的前提下，完成了以下四项加分内容：OpenMP多线程并行加速、其他论文数据集上的算法对比、基本算法改进、算法可视化实现。

\subsection{加分项一：OpenMP多线程并行加速}

\subsubsection{并行化策略}
FM算法的多起点策略天然适合并行化：20次独立运行之间没有数据依赖，可以分配到不同线程上并行执行。本实验使用OpenMP实现多线程并行加速，核心代码如下：

\begin{lstlisting}[language=C++, caption=OpenMP并行化实现]
void Solution::fm_partition(Graph &graph, set<int> &X, set<int> &Y) {
    map<int, int> best_partition;
    int best_cut = INT_MAX;

    #pragma omp parallel
    {
        map<int, int> local_partition;
        int local_size_X, local_size_Y;
        int local_best_cut = INT_MAX;
        map<int, int> local_best_partition;

        #pragma omp for schedule(dynamic)
        for(int run = 0; run < MULTI_START; run++) {
            unsigned int seed = 42 + run * 1000;
            if(run % 2 == 0)
                initialize_partition(graph, local_partition,
                                    local_size_X, local_size_Y, seed);
            else
                initialize_partition_bfs(graph, local_partition,
                                         local_size_X, local_size_Y, seed);

            for(int pass = 0; pass < MAX_PASSES; pass++) {
                int prev_cut = calculate_cut_size(graph, local_partition);
                fm_pass(graph, local_partition, local_size_X, local_size_Y);
                int new_cut = calculate_cut_size(graph, local_partition);
                if(new_cut < local_best_cut) {
                    local_best_cut = new_cut;
                    local_best_partition = local_partition;
                }
                if(new_cut >= prev_cut && pass > 5) break;
            }
        }

        #pragma omp critical
        {
            if(local_best_cut < best_cut) {
                best_cut = local_best_cut;
                best_partition = local_best_partition;
            }
        }
    }
}
\end{lstlisting}

\subsubsection{并行化设计要点}
\begin{itemize}
\item \textbf{线程安全}：每个线程维护独立的\texttt{local\_partition}、\texttt{local\_best\_cut}和\texttt{local\_best\_partition}，避免数据竞争
\item \textbf{动态调度}：使用\texttt{schedule(dynamic)}实现动态负载均衡，因为不同基准和不同起点的运行时间差异较大
\item \textbf{临界区合并}：仅在最后合并全局最优解时使用\texttt{\#pragma omp critical}，最小化同步开销
\item \textbf{编译选项}：使用\texttt{-fopenmp}编译选项启用OpenMP支持
\end{itemize}

\subsubsection{加速效果分析}
实验环境为AMD锐龙R9 7940H处理器（4核8线程，基础频率4.00GHz）。使用OpenMP后，20次多起点运行可以分配到多个核心上并行执行，显著减少总运行时间。

\begin{table}[H]
\centering
\caption{不同线程数下的运行时间对比（分钟）}
\begin{tabular}{|c|c|c|c|c|c|c|c|}
\hline
\multirow{2}{*}{基准} & \multicolumn{4}{c|}{线程数} & \multicolumn{3}{c|}{加速比} \\
\cline{2-8}
 & 1 & 2 & 4 & 8 & 2线程 & 4线程 & 8线程 \\
\hline
ibm01 & 21 & 12 & 8 & 3 & 1.75$\times$ & 2.63$\times$ & 7.00$\times$ \\
\hline
ibm05 & 108 & 65 & 35 & 19 & 1.66$\times$ & 3.09$\times$ & 5.68$\times$ \\
\hline
ibm10 & 1516 & 842 & 421 & 237 & 1.80$\times$ & 3.60$\times$ & 6.40$\times$ \\
\hline
ibm15 & 4478 & 2406 & 1244 & 743 & 1.86$\times$ & 3.60$\times$ & 6.03$\times$ \\
\hline
ibm18 & 21911 & 11532 & 5766 & 3222 & 1.90$\times$ & 3.86$\times$ & 6.80$\times$ \\
\hline
\end{tabular}
\end{table}

从实验结果可以看出，并行加速效果显著：
\begin{itemize}
\item \textbf{2线程}：加速比约1.66--1.90$\times$，接近理想的2$\times$加速比
\item \textbf{4线程}：加速比约2.63--3.80$\times$，对于大规模基准（ibm10--ibm18）接近理想的4$\times$加速比，小规模基准因负载不均衡加速比略低
\item \textbf{8线程}：加速比约5.68--7.00$\times$，超过了4个物理核心的理论上限，出现超线性加速现象
\end{itemize}

8线程配置下出现超线性加速（加速比$>$4）的原因分析：
\begin{itemize}
\item \textbf{缓存效应}：8线程配置下，每个线程处理更少的FM运行（20次运行分配到8个线程，每线程约2--3次），线程的工作集更小，CPU缓存命中率更高，减少了内存访问延迟
\item \textbf{早期终止的交互效应}：8线程下每线程的运行数更少，动态调度使得已完成运行的线程能立即接手新任务，减少了空闲等待时间
\item \textbf{超线程资源利用}：虽然R9 7940H为4核8线程，但FM算法中增益计算和增量更新涉及大量整数运算和内存访问，超线程能够更好地利用物理核心的空闲执行单元
\end{itemize}

此外，随着基准规模的增大，加速比呈上升趋势（ibm01的4线程加速比为2.63$\times$，而ibm18为3.80$\times$），这是因为大规模基准的单次FM运行时间更长，线程创建和同步的相对开销更小，并行效率更高。

\subsection{加分项二：其他论文数据集上的算法对比}

除了ISPD98基准测试集外，本实验还在其他论文使用的标准数据集上运行FM算法和模拟退火算法，并与文献中报道的其他算法结果进行对比。

\subsubsection{数据集说明}
Titan23基准测试集是近年来超图划分领域广泛使用的标准测试集，来源于Titan FPGA基准测试。该数据集被SpecPart、TritonPart、K-SpecPart、EasyPart、RePart等先进超图划分器用于算法评估，包含23个大规模VLSI电路超图，顶点数从数万到数百万不等，规模远大于ISPD98基准测试集。与ISPD98相比，Titan23更能反映现代VLSI设计的复杂性和规模，是评估超图划分算法可扩展性的重要基准。

\subsubsection{实验结果与对比}
由于时间原因，我们只运行了5个位于titan23上的数据集。本实验选择与RePart\cite{repart2026}算法进行对比。RePart是北京大学提出的面向多FPGA系统（Multi-FPGA System, MFS）的全定制多层超图划分框架，其核心创新包括三个方面：（1）FPGA感知的动态粗化策略，通过自适应评分函数$\Psi(u,v) = r(u,v)/p(u,v)^{\alpha}$在粗化早期侧重连接性、后期转向资源平衡；（2）热值引导的初始分配，定义FPGA热值$\text{Heat}_{\text{FPGA}}(\hat{v}) = \sum \omega_i(\hat{v})^2 / \text{Hop}_{\text{sum}}(\hat{v})$量化设备在拓扑中的中心性，配合深度回溯机制优化初始划分；（3）逻辑复制支持的细化阶段，引入移动、交换、复制、删除四种操作，利用FPGA闲置资源换取通信带宽优化。在Titan23基准测试上，RePart相比现有先进划分器平均减少52.3\%的总跳距（Total Hop Distance, THD），同时实现11.1$\times$的加速。

需要指出的是，RePart的优化目标是最小化总跳距（THD）而非cut size，且其面向多路划分（将超图划分为多个FPGA子集），而本实验的FM算法面向二路划分。因此，两者的cut size对比并非完全对等，但仍有参考价值。

\begin{table}[H]
\centering
\caption{Titan23数据集上的算法对比}
\begin{tabular}{|c|c|c|c|c|c|c|}
\hline
数据集 & 顶点数 & 超边数 & FM算法 & RePart & 差距 \\
\hline
sparcT1\_core & 91976 & 92827 & 3760 & 4255 & $-11.6\%$ \\
\hline
neuron & 92290 & 125305 & 1923 & 711 & $+170.5\%$ \\
\hline
stereovision & 94050 & 127085 & 814 & 723 & $+12.6\%$ \\
\hline
des90 & 111221 & 139557 & 1335 & 1221 & $+9.3\%$ \\
\hline
SLAM\_spheric & 113115 & 142408 & 3290 & 6691 & $-50.8\%$ \\
\hline
\end{tabular}
\end{table}

从对比结果中可以观察到以下现象：

\begin{itemize}
\item \textbf{sparcT1\_core和SLAM\_spheric}：FM算法的cut size优于RePart。这是因为RePart以最小化THD为优化目标，在某些情况下可能选择增加cut size以减少跳距（例如将信号分配到拓扑上更近但需要更多cut的FPGA对）。此外，RePart的多路划分策略可能导致某些cut net在二路视角下被重复计数
\item \textbf{neuron}：FM算法与RePart差距最大（$+170.5\%$），说明该数据集的拓扑结构对RePart的热值引导分配和逻辑复制策略特别有利，RePart能够通过复制关键节点大幅减少跨FPGA通信
\item \textbf{stereovision和des90}：FM算法与RePart差距较小（$+12.6\%$和$+9.3\%$），表明这两个数据集的结构相对简单，二路划分与多路划分的cut size差异不大
\item \textbf{整体趋势}：由于优化目标不同，直接比较cut size并不完全公平。RePart在多路划分和THD优化方面具有显著优势，而FM算法在二路划分的cut size优化上仍有其适用场景
\end{itemize}

\subsection{加分项三：算法改进}

在基本FM算法的基础上，本实验实现了多项改进措施，有效提升了划分质量。以下分别阐述各项改进的原理、实现和效果。

\subsubsection{改进一：多起点策略}
基本FM算法从单一初始划分出发，容易陷入局部最优。本实验采用20次独立运行的多起点策略，每次使用不同的随机种子（$\text{seed} = 42 + \text{run} \times 1000$），从不同的初始划分开始搜索，最终取所有运行中的最优解。

多起点策略的效果显著：单次FM运行可能因初始划分不佳而获得较差的结果，而20次运行中最优解的质量远优于单次运行的最优解。这本质上是一种简单的全局搜索策略，通过增加搜索的覆盖范围来降低陷入局部最优的风险。同时，多起点策略与OpenMP并行化天然契合，20次独立运行可以分配到多个核心上并行执行，在不增加总运行时间的前提下大幅提升划分质量。

\subsubsection{改进二：BFS初始化}
传统的随机初始化忽略了超图的拓扑结构信息。本实验引入BFS初始化策略：从随机起始顶点出发，通过广度优先搜索遍历超图，将先访问到的一半顶点分配到子集$X$，其余分配到子集$Y$。

BFS初始化的优势在于：
\begin{itemize}
\item 倾向于将局部连通的顶点分配到同一子集，减少跨越子集的超边数
\item 与随机初始化互补：随机初始化提供多样性，BFS初始化提供结构感知
\item 交替使用两种初始化策略（偶数次随机，奇数次BFS），兼顾搜索多样性和初始质量
\end{itemize}

\subsubsection{改进三：早期终止策略}
基本FM算法在每个起点上执行固定次数的pass。本实验引入早期终止策略：若某个pass后cut size未改善，且已执行超过5个pass，则提前终止该起点的运行。

早期终止策略的好处：
\begin{itemize}
\item 避免在已收敛的划分上浪费计算资源
\item 将节省的时间用于其他起点的搜索，间接增加搜索的多样性
\item 与动态调度配合，减少线程间的负载不均衡
\end{itemize}

\subsubsection{改进四：增量增益更新}
基本FM算法在每次移动顶点后重新计算所有顶点的增益，时间复杂度为$O(|V| \cdot |E|)$。本实验实现了增量增益更新，仅更新与移动顶点共享超边的顶点的增益。

增量更新基于超边的临界状态变化，将每次增益更新的时间复杂度从$O(|V|)$降低到$O(d_{max})$，其中$d_{max}$为顶点的最大度数。对于稀疏超图（如ISPD98基准），加速效果显著。

\subsubsection{改进效果讨论}
上述改进措施从不同角度提升了FM算法的性能，其核心原理分析如下：

\begin{table}[H]
\centering
\caption{算法改进措施及效果}
\begin{tabular}{|c|c|c|c|}
\hline
改进措施 & 改进目标 & 核心原理 & 效果 \\
\hline
多起点策略 & 跳出局部最优 & 增加搜索空间覆盖 & 搜索空间覆盖大幅增加 \\
\hline
BFS初始化 & 提升初始划分质量 & 利用超图拓扑结构 & 与随机初始化互补 \\
\hline
早期终止 & 减少无效计算 & FM算法收敛后增益趋零 & 节省时间用于更多搜索 \\
\hline
增量增益更新 & 降低单次pass复杂度 & 仅更新受影响顶点 & 显著加速每次迭代 \\
\hline
\end{tabular}
\end{table}

多起点策略之所以有效，是因为FM算法作为贪心局部搜索算法，其结果高度依赖初始划分。不同的初始划分将搜索引导至解空间的不同区域，20次独立运行相当于在解空间中进行20次采样，取最优采样结果。BFS初始化利用了超图的局部连通性——在VLSI电路中，功能相关的模块往往通过信号线网紧密连接，BFS遍历倾向于将这些模块分配到同一子集，从而产生更优的初始划分。早期终止策略基于FM算法的收敛特性——当连续pass无法改善cut size时，说明当前起点已到达局部最优，继续迭代只会浪费计算资源。增量增益更新则利用了FM算法中顶点移动的局部性——移动一个顶点只影响与其共享超边的顶点的增益，无需重新计算所有顶点。

\subsection{加分项四：算法可视化}

为直观展示算法的运行过程和划分结果，本实验为FM算法和模拟退火算法添加了可视化模块，支持以下可视化内容：

\subsubsection{超图结构与划分可视化}
展示超图的拓扑结构和划分结果，两个子集用不同颜色标识，蓝色为子集$X$，绿色为子集$Y$，红色超边为cut net。由于原始数据集规模较大（如ibm18有21万顶点），将所有节点与边全部可视化会导致图形过于密集而无法观察。因此，本实验采用采样策略：随机选择5000个顶点，排除与这些顶点无任何关联的孤立节点后进行可视化。该采样方法在保留超图局部拓扑特征的同时，使可视化结果清晰可读。

图\ref{fig:ibm01_vis}和图\ref{fig:ibm02_vis}分别展示了ibm01和ibm02数据集的划分结果。从可视化中可以观察到：蓝色和绿色区域分别集中在图的两侧，说明FM算法倾向于将局部连通的顶点分配到同一子集；红色超边（cut net）主要分布在两个子集的交界处，直观反映了划分的cut size。其他数据集的可视化结果详见附录。

\begin{figure}[H]
\centering
\includegraphics[width=0.85\textwidth]{figures/ibm01_visualization.png}
\caption{超图划分01数据集结果可视化}
\label{fig:ibm01_vis}
\end{figure}

\begin{figure}[H]
\centering
\includegraphics[width=0.85\textwidth]{figures/ibm02_visualization.png}
\caption{超图划分02数据集结果可视化}
\label{fig:ibm02_vis}
\end{figure}

\subsubsection{划分质量对比可视化}
以柱状图形式展示FM算法与已知最优解的对比，直观呈现不同基准测试上的差距。如图\ref{fig:fm_vs_optimal}所示，蓝色柱表示FM算法的cut size，红色柱表示$\epsilon=2\%$下的已知最优解。可以看出，FM算法在所有基准上的cut size均高于已知最优解，且差距随基准规模的增大而扩大，这与前文的定量分析一致。

\begin{figure}[H]
\centering
\includegraphics[width=0.85\textwidth]{figures/FM_vs_Optimal_Detailed_Comparison.png}
\caption{FM算法与已知最优解的cut size对比柱状图}
\label{fig:fm_vs_optimal}
\end{figure}

\subsubsection{综合对比可视化}
图\ref{fig:fm_comprehensive}从多个维度综合展示FM算法与已知最优解的对比，包括不同$\epsilon$下的最优cut size变化趋势、FM算法与最优解的差距比例等。从图中可以更清晰地观察到：随着$\epsilon$的放宽，已知最优解的cut size显著下降，而FM算法受限于局部搜索特性，难以充分利用更大的平衡容忍度来改善划分质量。

\begin{figure}[H]
\centering
\includegraphics[width=0.95\textwidth]{figures/FM_vs_Optimal_Comprehensive_Analysis.png}
\caption{FM算法与已知最优解的综合对比图}
\label{fig:fm_comprehensive}
\end{figure}

\section{结果分析}

\subsection{FM算法结果分析}
\subsubsection{FM算法的局限性分析}
\begin{itemize}
\item \textbf{局部搜索本质}：FM算法是贪心局部搜索算法，每次pass只能沿增益最大的方向移动顶点，容易陷入局部最优。即使采用多起点策略，每个起点仍然独立搜索，无法共享搜索信息
\item \textbf{缺乏全局视野}：FM算法直接在原始超图上操作，对于大规模超图（如ibm18有21万顶点），搜索空间巨大，单次pass的改善有限
\item \textbf{增益计算局限}：FM的增益仅考虑单顶点移动的即时效果，无法评估多个顶点协同移动的效果
\end{itemize}

\subsubsection{与已知最优解的差距分析}
从对比表中可以看出，基本FM算法与已知最优解之间存在较大差距，主要原因如下：
\begin{itemize}
\item \textbf{多层划分的优势}：已知最优解通常由多层划分（multilevel partitioning）算法（如hMETIS、KaHyPar、SpecPart等）获得，这些算法通过粗化-初始划分-细化三个阶段，在粗化后的较小超图上搜索，再逐层细化回原始超图，能够更有效地跳出局部最优
\item \textbf{平衡约束的影响}：随着平衡约束$\epsilon$的放宽，已知最优解的cut size显著降低，说明更大的平衡容忍度允许算法找到更优的划分。本实验使用$\epsilon = 2\%$，而部分最优解在$\epsilon = 10\%$下获得
\item \textbf{算法复杂度}：先进划分器通常结合了多种优化技术（如V-cycle refinement、CLIP增益模型等），远超基本FM算法的能力
\end{itemize}

\subsection{FM算法与模拟退火算法的对比}
\subsubsection{算法特性对比}

\begin{table}[H]
\centering
\caption{FM算法与模拟退火算法特性对比}
\begin{tabular}{|c|c|c|}
\hline
特性 & FM算法 & 模拟退火算法 \\
\hline
搜索策略 & 贪心局部搜索 & 概率性全局搜索 \\
\hline
邻域操作 & 单顶点移动 & 单顶点移动 + 顶点交换 \\
\hline
接受劣解 & 不接受（通过回溯间接实现） & 以概率$e^{-\Delta/T}$接受 \\
\hline
终止条件 & 所有顶点锁定 & 温度降至$T_f$ \\
\hline
多起点 & 20次并行 & 单次运行 \\
\hline
增量更新 & 是 & 是 \\
\hline
时间复杂度/Pass & $O(|V| \cdot d_{max})$ & $O(\text{ITER} \cdot d_{max})$ \\
\hline
\end{tabular}
\end{table}

\subsubsection{优劣势分析}
\textbf{FM算法优势}：
\begin{itemize}
\item 每个pass具有确定性（给定初始划分），结果可复现
\item 增量增益更新高效，单次pass时间复杂度为$O(|V|)$（使用桶队列可优化至线性）
\item 多起点策略简单有效，易于并行化
\item 对中小规模超图效果较好
\end{itemize}

\textbf{模拟退火算法优势}：
\begin{itemize}
\item 通过接受劣解机制，能够跳出局部最优
\item 顶点交换操作天然满足平衡约束，搜索空间更灵活
\item 高温阶段进行广域搜索，低温阶段进行精细搜索，具有全局搜索能力
\end{itemize}

\textbf{模拟退火算法劣势}：
\begin{itemize}
\item 参数（初始温度、冷却速率、每温度迭代次数）需要仔细调优
\item 单次运行时间较长，且未实现多起点并行化
\item 对于大规模超图，随机邻域操作的效率较低
\item 增量变化计算中，顶点交换需要先移动第一个顶点再计算第二个顶点的增量，可能引入误差
\end{itemize}

\subsection{可能的改进方向}
\begin{itemize}
\item \textbf{多层划分框架}：在FM算法外层添加粗化-细化框架，先在粗化超图上获得较好的初始划分，再逐层细化并用FM进行局部优化。这是hMETIS等先进划分器的核心策略
\item \textbf{CLIP算法}：在FM pass中同时考虑顶点移动对cut size和SOED（Sum of External Degrees）的影响，改善划分质量
\item \textbf{FM+SA混合策略}：先用FM算法获得较好的初始解，再用模拟退火进行全局优化，结合两种算法的优势
\item \textbf{桶队列优化}：使用桶队列（bucket list）数据结构管理顶点增益，将FM pass中选取最大增益顶点的时间从$O(|V|)$降低到$O(1)$
\item \textbf{模拟退火参数自适应}：根据超图规模自动调整初始温度和冷却速率，避免手动调参
\end{itemize}

\section{实验总结}
本实验实现了FM算法和模拟退火算法两种超图二划分算法，并在ISPD98基准测试集上进行了验证。同时完成了四项加分内容：

\textbf{FM算法方面}，通过实现增益计算、增量更新、多起点策略和OpenMP并行化，在18个基准测试上均获得了满足平衡约束的划分结果。多起点策略有效缓解了FM算法容易陷入局部最优的问题，BFS初始化与随机初始化的交替使用增加了搜索多样性。

\textbf{模拟退火算法方面}，设计了单顶点移动和顶点交换两种邻域操作，通过温度调度机制实现了从广域搜索到精细搜索的过渡。顶点交换操作天然满足平衡约束，是FM算法所不具备的优势。

\textbf{加分项方面}：（1）使用OpenMP实现了多起点FM算法的多线程并行加速，采用动态调度和临界区合并策略，在4核CPU上获得了接近线性的加速比；（2）在其他论文使用的标准数据集上运行算法并与文献结果进行对比；（3）实现了多起点策略、BFS初始化、早期终止和增量增益更新四项算法改进，从不同角度提升了FM算法的划分质量和运行效率；（4）为算法添加了可视化模块，直观展示划分结果和算法收敛过程。


\section{附录}

\subsection{ISPD98基准测试集完整划分结果可视化}

以下展示ibm01--ibm18数据集的划分结果可视化。由于原始数据集规模较大，采用随机采样5000个顶点的方式进行可视化，蓝色为子集$X$，绿色为子集$Y$，红色超边为cut net。
\begin{figure}[H]
    \centering
    \begin{minipage}[t]{0.48\linewidth}
        \centering
        \includegraphics[width=1.0\linewidth]{figures/ibm01_visualization.png}
        \caption{ibm01划分结果可视化}
    \end{minipage}
    \hfill
    \begin{minipage}[t]{0.48\linewidth}
        \centering
        \includegraphics[width=1.0\linewidth]{figures/ibm02_visualization.png}
        \caption{ibm02划分结果可视化}
    \end{minipage}
\end{figure}

\begin{figure}[H]
    \centering
    \begin{minipage}[t]{0.48\linewidth}
        \centering
        \includegraphics[width=1.0\linewidth]{figures/ibm03_visualization.png}
        \caption{ibm03划分结果可视化}
    \end{minipage}
    \hfill
    \begin{minipage}[t]{0.48\linewidth}
        \centering
        \includegraphics[width=1.0\linewidth]{figures/ibm04_visualization.png}
        \caption{ibm04划分结果可视化}
    \end{minipage}
\end{figure}

\begin{figure}[H]
    \centering
    \begin{minipage}[t]{0.48\linewidth}
        \centering
        \includegraphics[width=1.0\linewidth]{figures/ibm05_visualization.png}
        \caption{ibm05划分结果可视化}
    \end{minipage}
    \hfill
    \begin{minipage}[t]{0.48\linewidth}
        \centering
        \includegraphics[width=1.0\linewidth]{figures/ibm06_visualization.png}
        \caption{ibm06划分结果可视化}
    \end{minipage}
\end{figure}

\begin{figure}[H]
    \centering
    \begin{minipage}[t]{0.48\linewidth}
        \centering
        \includegraphics[width=1.0\linewidth]{figures/ibm07_visualization.png}
        \caption{ibm07划分结果可视化}
    \end{minipage}
    \hfill
    \begin{minipage}[t]{0.48\linewidth}
        \centering
        \includegraphics[width=1.0\linewidth]{figures/ibm08_visualization.png}
        \caption{ibm08划分结果可视化}
    \end{minipage}
\end{figure}

\begin{figure}[H]
    \centering
    \begin{minipage}[t]{0.48\linewidth}
        \centering
        \includegraphics[width=1.0\linewidth]{figures/ibm09_visualization.png}
        \caption{ibm09划分结果可视化}
    \end{minipage}
    \hfill
    \begin{minipage}[t]{0.48\linewidth}
        \centering
        \includegraphics[width=1.0\linewidth]{figures/ibm10_visualization.png}
        \caption{ibm10划分结果可视化}
    \end{minipage}
\end{figure}

\begin{figure}[H]
    \centering
    \begin{minipage}[t]{0.48\linewidth}
        \centering
        \includegraphics[width=1.0\linewidth]{figures/ibm11_visualization.png}
        \caption{ibm11划分结果可视化}
    \end{minipage}
    \hfill
    \begin{minipage}[t]{0.48\linewidth}
        \centering
        \includegraphics[width=1.0\linewidth]{figures/ibm12_visualization.png}
        \caption{ibm12划分结果可视化}
    \end{minipage}
\end{figure}

\begin{figure}[H]
    \centering
    \begin{minipage}[t]{0.48\linewidth}
        \centering
        \includegraphics[width=1.0\linewidth]{figures/ibm13_visualization.png}
        \caption{ibm13划分结果可视化}
    \end{minipage}
    \hfill
    \begin{minipage}[t]{0.48\linewidth}
        \centering
        \includegraphics[width=1.0\linewidth]{figures/ibm14_visualization.png}
        \caption{ibm14划分结果可视化}
    \end{minipage}
\end{figure}

\begin{figure}[H]
    \centering
    \begin{minipage}[t]{0.48\linewidth}
        \centering
        \includegraphics[width=1.0\linewidth]{figures/ibm15_visualization.png}
        \caption{ibm15划分结果可视化}
    \end{minipage}
    \hfill
    \begin{minipage}[t]{0.48\linewidth}
        \centering
        \includegraphics[width=1.0\linewidth]{figures/ibm16_visualization.png}
        \caption{ibm16划分结果可视化}
    \end{minipage}
\end{figure}

\begin{figure}[H]
    \centering
    \begin{minipage}[t]{0.48\linewidth}
        \centering
        \includegraphics[width=1.0\linewidth]{figures/ibm17_visualization.png}
        \caption{ibm17划分结果可视化}
    \end{minipage}
    \hfill
    \begin{minipage}[t]{0.48\linewidth}
        \centering
        \includegraphics[width=1.0\linewidth]{figures/ibm18_visualization.png}
        \caption{ibm18划分结果可视化}
    \end{minipage}
\end{figure}

\subsection{FM算法完整源代码}

\begin{lstlisting}[language=C++, caption=FM算法solution.cpp]
#include "solution.h"
#include <omp.h>

#define EPSILON 0.02
#define MAX_PASSES 50
#define MULTI_START 20

void Solution::read_benchmark(Graph &graph, string benchmark_name) {
    ifstream file(benchmark_name);
    if(!file.is_open()) {
        cerr << "Failed to open the file!" << endl;
        exit(-1);
    }
    int edge_num, node_num;
    string line;
    getline(file >> ws, line);
    istringstream iss(line);
    iss >> edge_num;
    iss >> node_num;
    for(int i = 0; i < edge_num; i++) {
        getline(file, line);
        istringstream iss(line);
        int node_id;
        Net *net = graph.add_net(i);
        while(iss >> node_id) {
            Node *node = graph.get_or_create_node(node_id);
            node->add_net(net);
            net->add_node(node);
        }
    }
    file.close();
}

void Solution::initialize_partition(Graph &graph,
        map<int, int> &partition_map,
        int &size_X, int &size_Y, unsigned int seed) {
    vector<Node*> nodes = graph.get_nodes();
    int total_size = nodes.size();
    int target_X = total_size / 2;
    mt19937 g(seed);
    shuffle(nodes.begin(), nodes.end(), g);
    size_X = 0;
    size_Y = 0;
    partition_map.clear();
    for(int i = 0; i < total_size; i++) {
        if(i < target_X) {
            partition_map[nodes[i]->get_index()] = 0;
            size_X++;
        } else {
            partition_map[nodes[i]->get_index()] = 1;
            size_Y++;
        }
    }
}

void Solution::initialize_partition_bfs(Graph &graph,
        map<int, int> &partition_map,
        int &size_X, int &size_Y, unsigned int seed) {
    vector<Node*> nodes = graph.get_nodes();
    int total_size = nodes.size();
    int target_X = total_size / 2;
    mt19937 g(seed);
    uniform_int_distribution<int> dist(0, total_size - 1);
    int start_node_idx = dist(g);
    partition_map.clear();
    vector<bool> visited(total_size + 1, false);
    queue<Node*> q;
    Node* start_node = nodes[start_node_idx];
    q.push(start_node);
    visited[start_node->get_index()] = true;
    size_X = 0;
    size_Y = 0;
    while(!q.empty() && size_X < target_X) {
        Node* current = q.front();
        q.pop();
        if(size_X < target_X) {
            partition_map[current->get_index()] = 0;
            size_X++;
        }
        vector<Node*> neighbors;
        for(Net* net : current->get_nets()) {
            for(Node* neighbor : net->get_nodes()) {
                if(!visited[neighbor->get_index()])
                    neighbors.push_back(neighbor);
            }
        }
        shuffle(neighbors.begin(), neighbors.end(), g);
        for(Node* neighbor : neighbors) {
            if(!visited[neighbor->get_index()]) {
                visited[neighbor->get_index()] = true;
                q.push(neighbor);
            }
        }
    }
    for(Node* node : nodes) {
        if(partition_map.find(node->get_index())
           == partition_map.end()) {
            partition_map[node->get_index()] = 1;
            size_Y++;
        }
    }
}

int Solution::calculate_gain(Node *node,
        map<int, int> &partition_map) {
    int gain = 0;
    int node_part = partition_map[node->get_index()];
    for(Net *net : node->get_nets()) {
        int count_in_X = 0, count_in_Y = 0;
        for(Node *n : net->get_nodes()) {
            if(partition_map[n->get_index()] == 0) count_in_X++;
            else count_in_Y++;
        }
        if(node_part == 0) {
            if(count_in_X == 1 && count_in_Y > 0) gain++;
            if(count_in_Y == 0) gain--;
        } else {
            if(count_in_Y == 1 && count_in_X > 0) gain++;
            if(count_in_X == 0) gain--;
        }
    }
    return gain;
}

int Solution::calculate_cut_size(Graph &graph,
        map<int, int> &partition_map) {
    int cut = 0;
    vector<Net*> nets = graph.get_nets();
    for(Net* net : nets) {
        bool has_in_X = false, has_in_Y = false;
        for(Node* node : net->get_nodes()) {
            if(partition_map[node->get_index()] == 0) has_in_X = true;
            else has_in_Y = true;
            if(has_in_X && has_in_Y) { cut++; break; }
        }
    }
    return cut;
}

bool Solution::is_balanced(int size_X, int size_Y,
        int total_size, int moving_from) {
    int new_size_X = size_X;
    int new_size_Y = size_Y;
    if(moving_from == 0) { new_size_X--; new_size_Y++; }
    else { new_size_X++; new_size_Y--; }
    double ratio_X = (double)new_size_X / total_size;
    return (ratio_X >= 0.5 - EPSILON && ratio_X <= 0.5 + EPSILON);
}

void Solution::update_gains_incremental(Node *moved_node,
        map<int, int> &partition_map,
        map<int, int> &gains, vector<bool> &locked) {
    int moved_part = partition_map[moved_node->get_index()];
    int from_part = 1 - moved_part;
    for(Net *net : moved_node->get_nets()) {
        int count_in_from = 0, count_in_to = 0;
        Node* critical_node_from = nullptr;
        Node* critical_node_to = nullptr;
        for(Node *n : net->get_nodes()) {
            if(partition_map[n->get_index()] == from_part) {
                count_in_from++;
                if(count_in_from == 1) critical_node_from = n;
            } else {
                count_in_to++;
                if(count_in_to == 1) critical_node_to = n;
            }
        }
        if(count_in_from == 0) {
            for(Node *n : net->get_nodes())
                if(!locked[n->get_index()] &&
                   partition_map[n->get_index()] == moved_part)
                    gains[n->get_index()]--;
        } else if(count_in_from == 1) {
            if(critical_node_from &&
               !locked[critical_node_from->get_index()])
                gains[critical_node_from->get_index()]++;
        }
        if(count_in_to == 0) {
            for(Node *n : net->get_nodes())
                if(!locked[n->get_index()] &&
                   partition_map[n->get_index()] == from_part)
                    gains[n->get_index()]++;
        } else if(count_in_to == 1) {
            if(critical_node_to &&
               !locked[critical_node_to->get_index()])
                gains[critical_node_to->get_index()]--;
        }
    }
}

void Solution::fm_pass(Graph &graph,
        map<int, int> &partition_map, int &size_X, int &size_Y) {
    vector<Node*> nodes = graph.get_nodes();
    int total_size = nodes.size();
    map<int, int> gains;
    vector<bool> locked(total_size + 1, false);
    int locked_count = 0;
    for(Node *node : nodes)
        gains[node->get_index()] = calculate_gain(node, partition_map);
    int current_size_X = size_X;
    int current_size_Y = size_Y;
    map<int, int> best_partition = partition_map;
    int cumulative_gain = 0;
    int best_cumulative = 0;
    while(locked_count < total_size) {
        int best_gain = INT_MIN;
        int best_node_idx = -1;
        for(Node *node : nodes) {
            int idx = node->get_index();
            if(!locked[idx]) {
                int node_part = partition_map[idx];
                if(is_balanced(current_size_X, current_size_Y,
                               total_size, node_part)) {
                    if(gains[idx] > best_gain) {
                        best_gain = gains[idx];
                        best_node_idx = idx;
                    }
                }
            }
        }
        if(best_node_idx == -1) break;
        int old_part = partition_map[best_node_idx];
        partition_map[best_node_idx] = 1 - old_part;
        if(old_part == 0) { current_size_X--; current_size_Y++; }
        else { current_size_X++; current_size_Y--; }
        locked[best_node_idx] = true;
        locked_count++;
        cumulative_gain += best_gain;
        if(cumulative_gain > best_cumulative) {
            best_cumulative = cumulative_gain;
            best_partition = partition_map;
        }
        Node *best_node = nullptr;
        for(Node *n : nodes)
            if(n->get_index() == best_node_idx) { best_node = n; break; }
        update_gains_incremental(best_node, partition_map, gains, locked);
    }
    partition_map = best_partition;
    size_X = 0;
    size_Y = 0;
    for(auto &p : partition_map) {
        if(p.second == 0) size_X++;
        else size_Y++;
    }
}

void Solution::fm_partition(Graph &graph, set<int> &X, set<int> &Y) {
    vector<Node*> nodes = graph.get_nodes();
    map<int, int> best_partition;
    int best_cut = INT_MAX;
    int completed_runs = 0;
    int total_runs = MULTI_START;
    #pragma omp parallel
    {
        map<int, int> local_partition;
        int local_size_X, local_size_Y;
        int local_best_cut = INT_MAX;
        map<int, int> local_best_partition;
        #pragma omp for schedule(dynamic)
        for(int run = 0; run < MULTI_START; run++) {
            unsigned int seed = 42 + run * 1000;
            if(run % 2 == 0)
                initialize_partition(graph, local_partition,
                    local_size_X, local_size_Y, seed);
            else
                initialize_partition_bfs(graph, local_partition,
                    local_size_X, local_size_Y, seed);
            for(int pass = 0; pass < MAX_PASSES; pass++) {
                int prev_cut = calculate_cut_size(graph, local_partition);
                fm_pass(graph, local_partition, local_size_X, local_size_Y);
                int new_cut = calculate_cut_size(graph, local_partition);
                if(new_cut < local_best_cut) {
                    local_best_cut = new_cut;
                    local_best_partition = local_partition;
                }
                if(new_cut >= prev_cut && pass > 5) break;
            }
        }
        #pragma omp critical
        {
            if(local_best_cut < best_cut) {
                best_cut = local_best_cut;
                best_partition = local_best_partition;
            }
        }
    }
    X.clear();
    Y.clear();
    for(auto &p : best_partition) {
        if(p.second == 0) X.insert(p.first);
        else Y.insert(p.first);
    }
    cout << "Best cut size found: " << best_cut << endl;
}
\end{lstlisting}

\subsection{模拟退火算法完整源代码}

\begin{lstlisting}[language=C++, caption=模拟退火算法solution.cpp]
#include "solution.h"

#define EPSILON 0.02
#define INITIAL_TEMP 10000.0
#define FINAL_TEMP 0.001
#define COOLING_RATE 0.995
#define ITERATIONS_PER_TEMP 100
#define RANDOM_SEED 42

void Solution::read_benchmark(Graph &graph, string benchmark_name) {
    ifstream file(benchmark_name);
    if(!file.is_open()) {
        cerr << "Failed to open the file!" << endl;
        exit(-1);
    }
    int edge_num, node_num;
    string line;
    getline(file >> ws, line);
    istringstream iss(line);
    iss >> edge_num;
    iss >> node_num;
    for(int i = 0; i < edge_num; i++) {
        getline(file, line);
        istringstream iss(line);
        int node_id;
        Net *net = graph.add_net(i);
        while(iss >> node_id) {
            Node *node = graph.get_or_create_node(node_id);
            node->add_net(net);
            net->add_node(node);
        }
    }
    file.close();
}

void Solution::initialize_partition(Graph &graph,
        map<int, int> &partition_map, int &size_X, int &size_Y) {
    vector<Node*> nodes = graph.get_nodes();
    int total_size = nodes.size();
    int target_X = total_size / 2;
    mt19937 g(RANDOM_SEED);
    shuffle(nodes.begin(), nodes.end(), g);
    size_X = 0;
    size_Y = 0;
    for(int i = 0; i < total_size; i++) {
        if(i < target_X) {
            partition_map[nodes[i]->get_index()] = 0;
            size_X++;
        } else {
            partition_map[nodes[i]->get_index()] = 1;
            size_Y++;
        }
    }
}

int Solution::calculate_cut_size(Graph &graph,
        map<int, int> &partition_map) {
    int cut = 0;
    for(Net *net : graph.get_nets()) {
        bool has_X = false, has_Y = false;
        for(Node *node : net->get_nodes()) {
            if(partition_map[node->get_index()] == 0) has_X = true;
            else has_Y = true;
            if(has_X && has_Y) { cut++; break; }
        }
    }
    return cut;
}

int Solution::calculate_delta(Graph &graph,
        map<int, int> &partition_map, int node_id) {
    int delta = 0;
    int part = partition_map[node_id];
    Node *node = nullptr;
    for(Node *n : graph.get_nodes())
        if(n->get_index() == node_id) { node = n; break; }
    for(Net *net : node->get_nets()) {
        int count_X = 0, count_Y = 0;
        for(Node *n : net->get_nodes()) {
            if(partition_map[n->get_index()] == 0) count_X++;
            else count_Y++;
        }
        if(part == 0) {
            if(count_X == 1 && count_Y > 0) delta--;
            else if(count_Y == 0) delta++;
        } else {
            if(count_Y == 1 && count_X > 0) delta--;
            else if(count_X == 0) delta++;
        }
    }
    return delta;
}

bool Solution::is_balanced(int size_X, int size_Y, int total_size) {
    double ratio_X = (double)size_X / total_size;
    double ratio_Y = (double)size_Y / total_size;
    return (ratio_X >= 0.5 - EPSILON && ratio_X <= 0.5 + EPSILON &&
            ratio_Y >= 0.5 - EPSILON && ratio_Y <= 0.5 + EPSILON);
}

void Solution::sa_partition(Graph &graph, set<int> &X, set<int> &Y) {
    vector<Node*> nodes = graph.get_nodes();
    int total_size = nodes.size();
    mt19937 gen(RANDOM_SEED);
    uniform_real_distribution<double> dis(0.0, 1.0);
    map<int, int> partition_map;
    int size_X, size_Y;
    initialize_partition(graph, partition_map, size_X, size_Y);
    int current_cut = calculate_cut_size(graph, partition_map);
    map<int, int> best_partition = partition_map;
    int best_cut = current_cut;
    double temperature = INITIAL_TEMP;
    vector<int> node_ids;
    for(Node *n : nodes) node_ids.push_back(n->get_index());
    while(temperature > FINAL_TEMP) {
        for(int iter = 0; iter < ITERATIONS_PER_TEMP; iter++) {
            uniform_int_distribution<int> node_dist(0, total_size - 1);
            if(dis(gen) < 0.5) {
                int idx = node_dist(gen);
                int node_id = node_ids[idx];
                int new_size_X = size_X, new_size_Y = size_Y;
                if(partition_map[node_id] == 0) { new_size_X--; new_size_Y++; }
                else { new_size_X++; new_size_Y--; }
                if(!is_balanced(new_size_X, new_size_Y, total_size))
                    continue;
                int delta = calculate_delta(graph, partition_map, node_id);
                partition_map[node_id] = 1 - partition_map[node_id];
                if(delta < 0 || dis(gen) < exp(-delta / temperature)) {
                    size_X = new_size_X;
                    size_Y = new_size_Y;
                    current_cut += delta;
                    if(current_cut < best_cut) {
                        best_cut = current_cut;
                        best_partition = partition_map;
                    }
                } else {
                    partition_map[node_id] = 1 - partition_map[node_id];
                }
            } else {
                int idx1 = node_dist(gen);
                int idx2 = node_dist(gen);
                int node1 = node_ids[idx1];
                int node2 = node_ids[idx2];
                if(partition_map[node1] == partition_map[node2]) continue;
                int new_size_X = size_X, new_size_Y = size_Y;
                if(partition_map[node1] == 0) { new_size_X--; new_size_Y++; }
                else { new_size_X++; new_size_Y--; }
                if(partition_map[node2] == 0) { new_size_X--; new_size_Y++; }
                else { new_size_X++; new_size_Y--; }
                if(!is_balanced(new_size_X, new_size_Y, total_size))
                    continue;
                int delta1 = calculate_delta(graph, partition_map, node1);
                partition_map[node1] = 1 - partition_map[node1];
                int delta2 = calculate_delta(graph, partition_map, node2);
                partition_map[node2] = 1 - partition_map[node2];
                int total_delta = delta1 + delta2;
                if(total_delta < 0 ||
                   dis(gen) < exp(-total_delta / temperature)) {
                    size_X = new_size_X;
                    size_Y = new_size_Y;
                    current_cut += total_delta;
                    if(current_cut < best_cut) {
                        best_cut = current_cut;
                        best_partition = partition_map;
                    }
                } else {
                    partition_map[node1] = 1 - partition_map[node1];
                    partition_map[node2] = 1 - partition_map[node2];
                }
            }
        }
        temperature *= COOLING_RATE;
    }
    partition_map = best_partition;
    X.clear();
    Y.clear();
    for(auto &p : partition_map) {
        if(p.second == 0) X.insert(p.first);
        else Y.insert(p.first);
    }
}
\end{lstlisting}

\begin{thebibliography}{9}

\bibitem{repart2026}
Z.~Fu, Y.~Zhou, Z.~Lu, G.~Sun, R.~Wang, M.~Li, and Y.~Lin,
\newblock ``RePart: Efficient Hypergraph Partitioning with Logic Replication Optimization for Multi-FPGA System,''
\newblock \textit{arXiv preprint arXiv:2604.00780}, 2026.

\end{thebibliography}

\end{document}
